\makeatletter
\@ifundefined{isChecklistMainFile}{
  % We are compiling a standalone document
  \newif\ifreproStandalone
  \reproStandalonetrue
}{
  % We are being \input into the main paper
  \newif\ifreproStandalone
  \reproStandalonefalse
}
\makeatother

\ifreproStandalone
\documentclass[letterpaper]{article}
\usepackage[submission]{aaai2027}
\setlength{\pdfpagewidth}{8.5in}
\setlength{\pdfpageheight}{11in}
\frenchspacing

\begin{document}
\fi
\setlength{\leftmargini}{20pt}
\makeatletter\def\@listi{\leftmargin\leftmargini \topsep .5em \parsep .5em \itemsep .5em}
\def\@listii{\leftmargin\leftmarginii \labelwidth\leftmarginii \advance\labelwidth-\labelsep \topsep .4em \parsep .4em \itemsep .4em}
\def\@listiii{\leftmargin\leftmarginiii \labelwidth\leftmarginiii \advance\labelwidth-\labelsep \topsep .4em \parsep .4em \itemsep .4em}\makeatother

\setcounter{secnumdepth}{0}
\renewcommand\thesubsection{\arabic{subsection}}
\renewcommand\labelenumi{\thesubsection.\arabic{enumi}}

\newcounter{checksubsection}
\newcounter{checkitem}[checksubsection]

\newcommand{\checksubsection}[1]{%
  \refstepcounter{checksubsection}%
  \paragraph{\arabic{checksubsection}. #1}%
  \setcounter{checkitem}{0}%
}

\newcommand{\checkitem}{%
  \refstepcounter{checkitem}%
  \item[\arabic{checksubsection}.\arabic{checkitem}.]%
}
\newcommand{\question}[2]{\normalcolor\checkitem #1 #2 \color{blue}}
\newcommand{\ifyespoints}[1]{\makebox[0pt][l]{\hspace{-15pt}\normalcolor #1}}

\section*{Reproducibility Checklist}

\vspace{1em}
\hrule
\vspace{1em}

\textbf{Instructions for Authors:}

This document outlines key aspects for assessing reproducibility. Please provide your input by editing this \texttt{.tex} file directly.

For each question (that applies), replace the ``Type your response here'' text with your answer.

\vspace{1em}
\noindent
\textbf{Example:} If a question appears as
%
\begin{center}
\noindent
\begin{minipage}{.9\linewidth}
\ttfamily\raggedright
\string\question \{Proofs of all novel claims are included\} \{(yes/partial/no)\} \\
Type your response here
\end{minipage}
\end{center}
you would change it to:
\begin{center}
\noindent
\begin{minipage}{.9\linewidth}
\ttfamily\raggedright
\string\question \{Proofs of all novel claims are included\} \{(yes/partial/no)\} \\
yes
\end{minipage}
\end{center}
%
Please make sure to:
\begin{itemize}\setlength{\itemsep}{.1em}
\item Replace ONLY the ``Type your response here'' text and nothing else.
\item Use one of the options listed for that question (e.g., \textbf{yes}, \textbf{no}, \textbf{partial}, or \textbf{NA}).
\item \textbf{Not} modify any other part of the \texttt{\string\question} command or any other lines in this document.\\
\end{itemize}

You can \texttt{\string\input} this .tex file right before \texttt{\string\end\{document\}} of your main file or compile it as a stand-alone document. Check the instructions on your conference's website to see if you will be asked to provide this checklist with your paper or separately.

\vspace{1em}
\hrule
\vspace{1em}

% The questions start here

\checksubsection{General Paper Structure}
\begin{itemize}

\question{Includes a conceptual outline and/or pseudocode description of AI methods introduced}{(yes/partial/no/NA)}
yes. The Conflict Witness Index subsection gives a conceptual description of
the maintained buckets, equivalence classes, constrained-value cells, update
path, query path, and shortest-certificate reconstruction.

\question{Clearly delineates statements that are opinions, hypothesis, and speculation from objective facts and results}{(yes/no)}
yes. Confirmatory synthetic results, the exploratory public-data topology
audit, formal guarantees, limitations, and future work are labeled separately.

\question{Provides well-marked pedagogical references for less-familiar readers to gain background necessary to replicate the paper}{(yes/no)}
partial. The paper cites the relevant primary literature on justifications,
diagnosis, provenance, incremental view maintenance, sparse and dense
retrieval, and rank fusion, but it is not a tutorial survey.

\end{itemize}
\checksubsection{Theoretical Contributions}
\begin{itemize}

\question{Does this paper make theoretical contributions?}{(yes/no)}
yes

	\ifyespoints{\vspace{1.2em}If yes, please address the following points:}
        \begin{itemize}
	
		\question{All assumptions and restrictions are stated clearly and formally}{(yes/partial/no)}
		yes. The store, schema, exact-literal, confidence, time/scope, query
		relevance, monotone context-cost, bounded-context, and rule-fragment
		assumptions are stated in the model section and collected again in the
		supplement.

		\question{All novel claims are stated formally (e.g., in theorem statements)}{(yes/partial/no)}
		yes. Conflict completeness is defined formally, and detector
		correctness, joint-budget feasibility and complexity, fixed-hop
		incompleteness, and the maintained-index contract are proposition
		statements.

		\question{Proofs of all novel claims are included}{(yes/partial/no)}
		yes. Core proofs are in the main paper and expanded proofs, including
		minimality and update-order details, are in the supplementary document.

		\question{Proof sketches or intuitions are given for complex and/or novel results}{(yes/partial/no)}
		yes. Each proposition is followed by a proof or proof sketch in the main
		paper, and the supplement provides the expanded argument.

		\question{Appropriate citations to theoretical tools used are given}{(yes/partial/no)}
		yes. Classical justification, diagnosis, provenance, and incremental
		view-maintenance results are cited where their concepts are used.

		\question{All theoretical claims are demonstrated empirically to hold}{(yes/partial/no/NA)}
		partial. The propositions are established by proof. Deterministic
		fixtures and implementation tests exercise their finite executable
		instances, but experiments do not establish the general theorems.

		\question{All experimental code used to eliminate or disprove claims is included}{(yes/no/NA)}
		yes. The anonymous code/data archive includes the negative controls,
		exact comparators, focused tests, and the corrected external-audit
		script and report.
	
	\end{itemize}
\end{itemize}

\checksubsection{Dataset Usage}
\begin{itemize}

\question{Does this paper rely on one or more datasets?}{(yes/no)}
yes

\ifyespoints{If yes, please address the following points:}
\begin{itemize}

		\question{A motivation is given for why the experiments are conducted on the selected datasets}{(yes/partial/no/NA)}
		yes. The synthetic fixtures isolate the retrieval contract under an
		oracle schema, while the separately labeled DBLP--Google Scholar audit
		tests only whether chained gold-match topology occurs outside the
		generator.

		\question{All novel datasets introduced in this paper are included in a data appendix}{(yes/partial/no/NA)}
		yes. All synthetic fixtures, instance manifests, expected labels, and
		frozen paper-facing results are included in the code/data archive.

		\question{All novel datasets introduced in this paper will be made publicly available upon publication of the paper with a license that allows free usage for research purposes}{(yes/partial/no/NA)}
		partial. The complete reviewer artifact is included with the submission;
		the public-release license remains an author approval item and is not
		presumed by the anonymous review package.

		\question{All datasets drawn from the existing literature (potentially including authors' own previously published work) are accompanied by appropriate citations}{(yes/no/NA)}
		yes. The DBLP--Google Scholar entity-resolution benchmark is cited in
		the paper.

		\question{All datasets drawn from the existing literature (potentially including authors' own previously published work) are publicly available}{(yes/partial/no/NA)}
		yes. The external audit uses the public CC-BY release identified by its
		citation and records the input hashes in the shipped report.

		\question{All datasets that are not publicly available are described in detail, with explanation why publicly available alternatives are not scientifically satisficing}{(yes/partial/no/NA)}
		NA. No non-public dataset supports a reported result.

\end{itemize}
\end{itemize}

\checksubsection{Computational Experiments}
\begin{itemize}

\question{Does this paper include computational experiments?}{(yes/no)}
yes

\ifyespoints{If yes, please address the following points:}
\begin{itemize}

		\question{This paper states the number and range of values tried per (hyper-) parameter during development of the paper, along with the criterion used for selecting the final parameter setting}{(yes/partial/no/NA)}
		partial. The final retrieval depths, top-$K$ grid, rank-fusion constant,
		confidence semantics, fixture families, joint-conflict counts, scale and
		class-stress sizes, and repetition counts are reported. The manuscript
		does not reconstruct every
		development-time trial.

		\question{Any code required for pre-processing data is included in the appendix}{(yes/partial/no)}
		yes. Synthetic data generation and the built-in-only external topology
		audit script are included in the code/data archive.

		\question{All source code required for conducting and analyzing the experiments is included in a code appendix}{(yes/partial/no)}
		partial. The complete synthetic-fixture, exact, maintained-index,
		joint-budget, class-stress, on-demand, sparse-retrieval, scoring, and
		scaling paths are included.
		The frozen dense outputs and pinned Python environment are included, but
		reproducing dense embeddings also requires obtaining the cited
		third-party model weights. For double-blind review, the project-owned
		RDF namespace is consistently replaced by a reserved
		\texttt{example.org} namespace without changing graph structure or
		scientific values.

		\question{All source code required for conducting and analyzing the experiments will be made publicly available upon publication of the paper with a license that allows free usage for research purposes}{(yes/partial/no)}
		partial. Source is available to reviewers now; the public-release license
		requires author approval before publication.

		\question{All source code implementing new methods have comments detailing the implementation, with references to the paper where each step comes from}{(yes/partial/no)}
		partial. The implementation is typed and tested and documents its
		invariants and certificate construction, but not every source step has
		an explicit paper-section cross-reference.

		\question{If an algorithm depends on randomness, then the method used for setting seeds is described in a way sufficient to allow replication of results}{(yes/partial/no/NA)}
		yes. Fixtures and outputs have deterministic identifiers or SHA-256
		digests, scale-query sampling is deterministic, and the artifact
		instructions distinguish exact regeneration from measured timing.

		\question{This paper specifies the computing infrastructure used for running experiments (hardware and software), including GPU/CPU models; amount of memory; operating system; names and versions of relevant software libraries and frameworks}{(yes/partial/no)}
		yes. The paper reports Node.js 22.17.1, Sentence Transformers 5.6.1,
		the pinned encoder revision, dependency locks, and the AMD Ryzen
		Threadripper PRO 5955WX timing host with 256 GiB RAM and AlmaLinux 9.7;
		no GPU was used.

		\question{This paper formally describes evaluation metrics used and explains the motivation for choosing these metrics}{(yes/partial/no)}
		yes. Conflict completeness and context cost are defined formally; the
		evaluation explains why complete minimal-witness inclusion, serialized
		triple count, joint witness-union cost, update amplification,
		materialized conflict pairs, and repeated timing are reported.

		\question{This paper states the number of algorithm runs used to compute each reported result}{(yes/no)}
		yes. Deterministic fixture and joint-budget results are exact single
		executions; scale and class-stress timings use one unrecorded warm-up and
		11 measured repetitions.

		\question{Analysis of experiments goes beyond single-dimensional summaries of performance (e.g., average; median) to include measures of variation, confidence, or other distributional information}{(yes/no)}
		yes. The paper reports per-family means and maxima, exact confusion
		counts, full family breakdowns, and timing medians with IQRs.

		\question{The significance of any improvement or decrease in performance is judged using appropriate statistical tests (e.g., Wilcoxon signed-rank)}{(yes/partial/no)}
		no. Deterministic completeness and context results are exact counts, and
		timing comparisons are descriptive medians and IQRs; the paper makes no
		statistical-significance claim.

		\question{This paper lists all final (hyper-)parameters used for each model/algorithm in the paper’s experiments}{(yes/partial/no/NA)}
		yes. The final retrieval depths, top-$K$ grid, rank-fusion constant,
		encoder revision, rule settings, joint-conflict counts, scale and stress
		sizes, query counts, and repetition protocol are in the paper and
		artifact.


\end{itemize}
\end{itemize}
\ifreproStandalone
\let\reproEndDocument\enddocument
\else
\let\reproEndDocument\relax
\fi
\reproEndDocument
